\begin{helper}
For every \(\eta>0\), every \(\varepsilon_1>0\), and every
\(\eta<1/16\), there is \(n_0\) such that for every \(n>n_0\), if
\(L=\log n\), then
\[
\frac{2(1+\eta)\sqrt L\,L}{n^{1/4-2\eta}}\le \frac14,
\]
\[
\frac{4(1+\eta)\sqrt L\,L}{n^\eta}\le \frac14,
\]
and
\[
\exp\!\left(-\frac12 n^{3/4-2\eta}\right)\le\varepsilon_1.
\]
\end{helper}
\begin{proof}
Put \(a=1/4-2\eta\) and \(b=3/4-2\eta\).  The hypothesis
\(\eta<1/16\) gives \(a>0\) and \(b>0\).  Since \(\log n=o(n^{a/2})\),
we may enlarge \(n_0\) so that
\[
L\le c_1 n^{a/2},\qquad c_1=\frac1{8(1+\eta)}.
\]
Also enlarge \(n_0\) so that \(L\ge1\).  Then \(\sqrt L\le L\), and hence
\[
\sqrt L\,L\le L^2\le c_1^2 n^a.
\]
The choice of \(c_1\) makes
\[
2(1+\eta)c_1^2\le \frac14,
\]
so
\[
2(1+\eta)\sqrt L\,L\le \frac14 n^a,
\]
which is equivalent to the first displayed inequality.

Similarly, because \(\log n=o(n^{\eta/2})\), after enlarging \(n_0\) we have
\[
L\le c_2 n^{\eta/2},\qquad c_2=\frac1{16(1+\eta)}.
\]
Again \(L\ge1\) gives
\[
\sqrt L\,L\le L^2\le c_2^2 n^\eta,
\]
and the choice of \(c_2\) gives
\[
4(1+\eta)c_2^2\le\frac14.
\]
This proves the second displayed inequality.

Finally, since \(b>0\), \(n^b\to\infty\).  Enlarge \(n_0\) so that
\[
n^b\ge \max\{0,-2\log\varepsilon_1\}.
\]
Then
\[
-\frac12n^b\le \log\varepsilon_1.
\]
Because \(\varepsilon_1>0\), exponentiating gives
\[
\exp\!\left(-\frac12 n^b\right)\le\varepsilon_1.
\]
\end{proof}
